\section{Stronger density estimates: multilinear bounds and a new prefix law}
\label{sec:density:moments-prefix}

\paragraph{What this section proves and how to read it.}
To obtain $\vartheta>1$ one needs a multilinear bound
$|\int\prod_vf_v\,d\pi|\le\prod_v\|f_v\|_{L^r}$ with $r<72$
(a \emph{unit mixed moment}, defined below; ``moment'' here means
this multilinear norm, not $\E X^k$). This section proves two
finite-alphabet inequalities giving such bounds and constructs a
replacement prefix law (independent slot labels from the same
clique) satisfying them. Inputs: \cref{sec:frontier:iidprefix}.
Outputs: the moment hypotheses used in \cref{sec:density-upgrade}.
Baseline results do not depend on this section.

\begin{definition}[Unit mixed moment]\label{def:density:mixed-moment}
Let $d\ge1$ and $r\ge1$. Let $\pi$ be a probability measure on $X^d$, with each coordinate uniform
on the finite set $X$. It has \emph{unit mixed $L^r$ norm} if
\[
 \left|\int\prod_{v=1}^d f_v(x_v)\,d\pi(x)\right|
 \le\prod_{v=1}^d\|f_v\|_{L^r(X)}
\]
for all real or complex functions $f_v$, where all norms use uniform
probability measure. Its \emph{density endpoint} is $d/r$. If $\pi'$ on $X^{d'}$ and $\pi''$ on $X^{d''}$ have unit mixed $L^r$ norm, so does $\pi'\otimes\pi''$, by successive integration; the same holds for convex mixtures and common coordinate permutations. A bound at $r$ also holds at every larger exponent.
\end{definition}

\subsection{Uniform convexity and mixtures on a finite alphabet}

The stronger induction needs moment bounds whose density endpoint $d/r$ exceeds $1$. The elementary inequality behind them is uniform
convexity of $L^r$ norms; we prove it, and the finite-alphabet
mixture estimate built on it, here because the exact constants
matter later.

\begin{lemma}[Uniform convexity]\label{lem:density:uniform-convexity}
On a finite probability space, for $1<r\le2$ and nonnegative $f$,
\[
 \|f\|_r^2\ge(\E f)^2+(r-1)\|f-\E f\|_r^2.
\]
\end{lemma}
\begin{proof}
For a positive function $g$ and a real direction $h$, differentiation gives
\begin{align*}
 \frac{d^2}{dt^2}\|g+th\|_r^2\bigg|_{t=0}
 &=2(r-1)\|g\|_r^{2-r}\E(g^{r-2}h^2)\\
 &\quad+2(2-r)\|g\|_r^{2-2r}\bigl(\E(g^{r-1}h)\bigr)^2
 \ge2(r-1)\|h\|_r^2.
\end{align*}
The last step follows from H\"older applied to
$|h|^r=(g^{r-2}h^2)^{r/2}g^{r(2-r)/2}$.
Take $g=\E f$, $h=f-\E f$, and integrate the second derivative along
$g+th$, $0\le t\le1$. The first derivative at zero vanishes. If $f$ has
zeros, apply the argument to $f+\varepsilon$ and let $\varepsilon$ decrease
to zero; the assertion is immediate when $\E f=0$.
\end{proof}

\begin{lemma}[Typed finite-alphabet mixture]\label{lem:density:typed-mixture}
Let $|X|=Q\ge2$ and $1\le s<R<d$. Put $\kappa=s/R$ and, for $R<v\le d$,
\[
 C_v(Q)=\frac d v Q^{(v/R-1)(d/v-1)}
 \frac{Q^{v/R}-(v/R)Q+v/R-1}
      {\kappa Q+1-\kappa-Q^\kappa}.
\]
Let $\pi=\sum_j w_j\pi_j$, where $w_j\ge0$, $\sum_jw_j=1$, and each
$\pi_j$ has unit mixed $L^{r_j}$ norm and uniform coordinate marginals.
Suppose $\sum_{j:r_j\le s}w_j\ge w_0$.
Components with $s<r_j\le R$ require no further condition. For every
remaining component choose $v_j\in(R,d]$ such that $\pi_j$ has unit mixed $L^{v_j}$ norm (for example $v_j=d$). If
\[
 \sum_{j:\,v_j>R}w_j C_{v_j}(Q)\le w_0,
\]
then $\pi$ has unit mixed $L^R$ norm. Moreover, if
$Q^{\kappa-1}\le\kappa/2$, then
\begin{equation}\label{eq:density:typed-mixture-coarse}
 C_v(Q)\le\frac{2d}{\kappa v}Q^{d/R-d/v}.
\end{equation}
In particular, a mixture of a unit mixed $L^s$ law of weight $1-t$ and
an arbitrary uniform-marginal law of weight $t$ has unit mixed $L^R$ norm
whenever $t(1+C_d(Q))\le1$.
\end{lemma}

\begin{proof}
We first prove a one-coordinate estimate. Let $f$ be a function on $X$ with $\|f\|_R=1$, and write $u=|f|^R$. Then $\E u=1$ and $0\le u\le Q$.
For $\lambda>1$ put
\[
 \phi(t)=t^\lambda-\lambda t+\lambda-1,
 \qquad\psi(t)=\kappa t+1-\kappa-t^\kappa.
\]
Both functions and their first derivatives vanish at $1$, and
$\phi''(t)/\psi''(t)$ is a positive constant times
$t^{\lambda-\kappa}$. It follows that
\[
 \frac{\phi(t)}{\psi(t)}\le\frac{\phi(Q)}{\psi(Q)}
 \quad(0\le t\le Q),
\]
with $\phi(1)/\psi(1)$ interpreted as $\phi''(1)/\psi''(1)=\lambda(\lambda-1)/(\kappa(1-\kappa))$. For completeness, on $t<1$
the ratio is a weighted average of $\phi''/\psi''$ on $[t,1]$, hence at most
its value at $1$. On $t>1$, first $\phi'/\psi'$ and then $\phi/\psi$ are
increasing weighted averages of the preceding increasing ratio.

Take $\lambda=v/R$ and $A=\E u^\lambda$. Since
$1\le A\le Q^{\lambda-1}$ and $d/v\ge1$,
\begin{align*}
 \|f\|_v^d-1
 &=A^{d/v}-1\\
 &\le\frac d v Q^{(\lambda-1)(d/v-1)}(A-1)\\
 &\le C_v(Q)(1-\E u^\kappa)
 \le C_v(Q)(1-\|f\|_s^d).
\end{align*}
The last inequality uses $\|f\|_s^d=(\E u^\kappa)^{d/s}\le\E u^\kappa$.
If $Q^{\kappa-1}\le\kappa/2$, then
$\psi(Q)\ge\kappa Q/2$ and $\phi(Q)\le Q^\lambda$; this proves
\eqref{eq:density:typed-mixture-coarse}.

For unequal inputs normalized by $\|f_v\|_R=1$, apply H\"older to the
component index with weights $w_j$:
\[
 |\Lambda_\pi((f_v)_v)|
 \le\sum_jw_j\prod_v\|f_v\|_{r_j}
 \le\prod_v\left(\sum_jw_j\|f_v\|_{r_j}^d\right)^{1/d}\le1.
\]
The final inequality follows from the preceding scalar bound, since the
weight $w_0$ supplies the deficit $1-\|f_v\|_s^d$. Ordinary H\"older gives
unit mixed $L^d$ norm for every law with uniform coordinate marginals,
which proves the final assertion.
\end{proof}

\subsection{A pair moment and independent draws from a digit clique}

The basic local input is one pair of labels whose difference is a
square, or a nonsquare. The next lemma gives its moment bound;
independent draws from the digit clique of the prefix are products
of such pairs, so the same bound applies to them.

\begin{lemma}[Square-pair moment]\label{lem:density:edge-moment}
Let $p$ be an odd prime. A uniform ordered strict square pair, or a
uniform ordered nonsquare pair, in $\Fp$ has unit mixed $L^{1+1/j}$ norm
whenever $j$ is a positive integer satisfying
\begin{equation}\label{eq:density:edge-integer-gate}
 j^{j+1}\left(\frac{p+1}{p-1}\right)^{j-1}
 \frac1{(\sqrt p-1)^2}\le1.
\end{equation}
The conclusion persists on $\Z/p^2\Z$ when each label is written as $x+py$ with $x,y\in\Fp$, the two $y$-coordinates are independent and uniform, and the pair of $x$-coordinates has either of these laws.
\end{lemma}
\begin{proof}
Relative to uniform probability measure, the centered bilinear form $(f,g)\mapsto\E fg-\E f\E g$ of the pair law has norm at most $L=(p+1)/(p-1)$ on $L^1\times L^1$ and at most $\rho=(\sqrt p+1)/(p-1)=1/(\sqrt p-1)$ on $L^2\times L^2$. The latter is the standard Paley
character-spectrum bound; it holds for $p\equiv1$ and $p\equiv3\pmod4$ and for either choice of square or nonsquare difference. Bilinear interpolation at $r=1+1/j$ gives centered
$L^r$ norm at most $\beta=L^{(j-1)/(j+1)}\rho^{2/(j+1)}$.
Under \eqref{eq:density:edge-integer-gate}, $\beta\le r-1$.
For nonnegative inputs of $L^r$ norm one and means $m,n$, the expectation of the product of the two inputs under the pair law is therefore at most
\[
 mn+\frac\beta{r-1}\sqrt{(1-m^2)(1-n^2)}\le1,
\]
by \cref{lem:density:uniform-convexity} and Cauchy--Schwarz in $\R^2$.
Positivity of the pair law allows absolute values of arbitrary inputs.
For the prime-square law, first average each input over its high digit;
this is an $L^r$ contraction.
\end{proof}

\begin{corollary}[Invariant distinct pairs]\label{lem:density:invariant-pair}
For $p\ge64$, any probability law on ordered distinct pairs in $\Fp$
which is invariant under simultaneous square-affine maps has unit mixed
$L^{3/2}$ norm.
\end{corollary}
\begin{proof}
Such a law is a convex mixture of the square and nonsquare pair laws:
these are the two orbits of the square-affine group on distinct pairs.
For $j=2$, the left side of \eqref{eq:density:edge-integer-gate} is at most
$8\cdot65/(63\cdot49)<1$ for $p\ge64$. Apply
\cref{lem:density:edge-moment} and mix.
\end{proof}

\begin{lemma}[Thirty-six digit-clique draws]\label{lem:density:clique36}
Let $2^{16}\le p<2^{512}$ be a prime congruent to $1$ modulo $4$.
Let $(A,B)\in(\Z/p^2\Z)^2$ have the digit-pair law of
\cref{def:clique-law}: its low digits are a uniform strict square pair,
and its high digits are independent and uniform. Conditional on $(A,B)$,
draw $36$ labels independently, each uniform on $\{A,B\}$. Their law has
unit mixed $L^{600/17}$ norm, equivalently density endpoint $51/50$.
\end{lemma}
\begin{proof}
Put $d=36$, $\theta=51/50$ and $R=d/\theta=600/17$. Choose the unique
$b\in\{16,32,\ldots,496\}$ with $2^b\le p<2^{b+16}$. In increasing order
of $b$, use respectively
\[
\begin{gathered}
 5,9,12,15,18,20,23,26,28,31,33,36,38,40,43,45,\\
 47,50,52,54,56,58,61,63,65,67,69,71,73,75,77
\end{gathered}
\]
for the integer $j$, and put $r=1+1/j$, $s=18r$, $\kappa=s/R$.
The condition \eqref{eq:density:edge-integer-gate} holds at $p=2^b$, with
$\sqrt p=2^{b/2}$, and hence for every $p$ in the dyadic
interval $[2^b,2^{b+16})$ (called a band only in this proof; not
an upper band) by monotonicity in $p$.

Expand the independent binary choices into their $2^{36}$ choice vectors.
The choice vectors with eighteen $A$'s and eighteen $B$'s have total weight $w_0=\binom{36}{18}/2^{36}$ and
unit mixed $L^s$ norm: apply the pair moment to the two products of
$18$ inputs and then H\"older within each product. For minority count
$0\le k<18$, the total weight of the two complementary classes is
$w_k=2\binom{36}{k}/2^{36}$. The corresponding law has unit mixed norm at
\[
 v_k=\begin{cases}
 36,&k=0,\\
 \min\{36,(36-k)r\},&1\le k<18.
 \end{cases}
\]
The truncation at $36$ (the $\min\{36,\cdot\}$ in the display;
not a spectral cap) follows from uniform marginals. Thus only the classes with $v_k>R$ occur in the sum in \cref{lem:density:typed-mixture}.

Here $Q=p^2$. The following finite rational tests hold on every band:
\begin{gather}
 2^{-\lfloor2b(1-\kappa)\rfloor}\le\kappa/2,
 \label{eq:density:clique-denominator}\\
 \frac1{w_0}\sum_{k:\,v_k>R}w_k\frac{72}{\kappa v_k}
 2^{\lceil2(b+16)(\theta-36/v_k)\rceil}
 \le\frac{95659840}{13884957009}<1.
 \label{eq:density:clique-binomial-gate}
\end{gather}
For each of the $31$ values of $b$ these are two explicit integer/rational inequalities: the first uses
$Q\ge2^{2b}$, and the second uses $Q\le2^{2(b+16)}$ and rounds every power
of two upward. The largest ratio occurs at $b=16$.
The standalone checker
\path{research/four_routes_2026_10_26/dense_moments/check_iid72_prefix_moments.py}
reproduces these tests, the listed interpolation gates and the price
bounds below using only integer and rational arithmetic; its output is
\path{research/four_routes_2026_10_26/dense_moments/iid72_prefix_moment_gates.json}.
Equations \eqref{eq:density:clique-denominator}--\eqref{eq:density:clique-binomial-gate}
and \cref{lem:density:typed-mixture} prove the assertion.
\end{proof}

\begin{lemma}[At most two copies per label]\label{lem:density:clique-copies}
Under the digit-pair law of \cref{lem:density:clique36}, fix a map from independent draws from $\{A,B\}$ to the $72$ physical slots under which each draw is used at most twice.
The resulting law has unit mixed $L^{1200/17}$ norm. The same holds after
tensoring $m$ independent digit pairs, for any number $m$ of
digit pairs, and applying any common random square-affine
permutation of $\Z/p^{2m}\Z$.
\end{lemma}
\begin{proof}
Write $C=\{A,B\}$ for the random two-element clique and $\nu_C$ for the uniform law on $C$.
Conditional Cauchy--Schwarz in each group of one or two slots gives
\[
 \E\left[\prod_v|f_v|\,\middle|\,C\right]
 \le\prod_v(\nu_C|f_v|^2)^{1/2}.
\]
H\"older on $C$ with $72$ equal exponents, followed by
\cref{lem:density:clique36} with identical input $|f_v|^2$, bounds the
unconditional expectation by
\[
 \prod_v\bigl[\E_C(\nu_C|f_v|^2)^{36}\bigr]^{1/72}
 \le\prod_v\|f_v\|_{1200/17}.
\]
The bound for one digit pair implies the bound for $m$ independent digit pairs, for arbitrary functions of the full digit vector, by successive integration. A fixed affine permutation
preserves the uniform norms, and its positive average preserves the bound.
\end{proof}

\subsection{The replacement prefix and its dominators}

With these moments the prefix law of \cref{sec:frontier:iidprefix}
can be strengthened: at large prefix primes the $72$ labels are
drawn independently from the square clique, instead of through six
chamber labels. The theorem states the properties of this
replacement law and of the laws dominating it, which remain valid
for density exponents up to $51/50$.

\begin{theorem}[A prefix with density endpoint $51/50$]
\label{thm:density:iid-prefix}
At each prefix prime $p\equiv1\pmod4$, keep the $m=\mcl=8192$
digit pairs, label group $\Z/p^{2m}\Z$, random square clique $\mathcal C_p$
and common square-affine map of \cref{def:clique-law}. For $p\ge2^{16}$,
replace its six-label pullback by $72$ conditionally independent uniform
draws from $\mathcal C_p$, one per slot. For $p<2^{16}$ retain
the law of \cref{def:clique-law}. Then:
\begin{enumerate}[label=(\alph*)]
\item The new reference is reflection positive at every cut of the slot array $V$ of \cref{def:slot-array} used in the construction,
 square-affine invariant, and has uniform one-slot marginals. Every
 nonzero difference of two labels in its support has the same square
 property as in \cref{lem:frontier:iidlocal}(iii).
\item At a large prefix prime, the reference, the diagonal copy of its
 actual reserved half, and every law obtained by identifying disjoint
 pairs among its independent labels have unit mixed $L^{1200/17}$ norm.
 These statements tensor over selected prefix atoms.
\item For the six actual queries, let $E_p$ be the restriction of the new reference law $\mu_p$ to the union of their equality events, and put
 \[
  \eta=1-(1-2^{-m})^6\le6\cdot2^{-m}.
 \]
 Both reserved-half marginals of $E_p$ equal $\eta h_p$, where $h_p$ is
 the actual $36$-label half marginal of the new reference $\mu_p$.
 With $D_p=\operatorname{Diag}_\iota(h_p)$ and $\nu_p=\mu_p-E_p$,
 \[
 \begin{pmatrix}\mu_p+\eta D_p&\nu_p\\
                 \nu_p^*&\mu_p+\eta D_p\end{pmatrix}\succeq_\iota0.
 \]
 The measure $\nu_p=\mu_p-E_p$ is nonnegative, has mass $1-\eta>0$, and accepts all six
 queries. The symmetric and skew parts of $E_p$ each have operator norm at most $\eta$ between the actual half $L^2$ spaces; no symmetry of $E_p$ is assumed.
\item For $0<a\le91/100000$ and $1<\theta\le51/50$, each of the three
 aggregate prefix allowances $S,A,T$ is less than $2\cdot10^{-13}$,
 including the factor $Q^{\theta-1}$, for alphabet size $Q$, at the unchanged small primes. Thus each is below the bound $10^{-12}$ assigned to it in \cref{def:closing-quantities}.
\end{enumerate}
\end{theorem}

\begin{proof}
Conditional on the clique and affine map, the $72$ new labels
are iid. At every reflection cut with two halves of $36$ slots
(a balanced reflection) the two halves are independent and have
the same law, so its cut form is the square of a conditional
mean.
This proves reflection positivity. The existing common affine averaging
gives invariance and uniform marginals. The digit argument for the square
property is unchanged: if the first digit pair at which two
clique elements differ is the pair in position $j$ (counted from
$0$), their difference has $p$-adic valuation $2j$ and unit part
congruent to $\pm(b_j-a_j)$ modulo $p$, and $-1$ is a
square. Hensel lifting then gives the required unit square. This proves
(a). Part (b) is \cref{lem:density:clique-copies}; the reserved diagonal
uses $36$ independent labels, each exactly twice.

The query endpoints on the original labels are pairwise
disjoint and each query crosses the reserved binary cut.
Conditional on the clique and on the labels of the plus half,
the six labels at the minus endpoints of the six queries are
independent uniform choices from a set of $2^m$ elements. Each equality has probability $2^{-m}$,
independently, so their union has probability $\eta$, independent of the
plus-half realization. The reverse argument gives the other marginal.
Every positive coupling is a contraction between its actual marginal
$L^2$ spaces. Hence
\[
 \begin{pmatrix}\eta D_p&E_p\\E_p^*&\eta D_p\end{pmatrix}
 \succeq_\iota0.
\]
Changing the signs of the off-diagonal blocks and adding the positive
block matrix with all four entries $\mu_p$ proves the displayed bracket
in (c). The accepted part removes precisely the zero query differences.
The symmetric and skew parts of its change have operator norms at most
$\eta$, and $E_p$ is dominated at $\iota$ by the diagonal envelope $\eta D_p$.
Using the separate symmetric and skew envelopes for matching instead
costs at most $2\eta D_p$. In all cases the chosen matching sides contain
$\mu_p$ plus a positive diagonal, so the lower bound for the mean constant $m$ from these sides is the same as for $\mu_p$.
The first-stage bracket at a prefix atom remains the one for $\mu_p$.

An individual query-equality measure has mass $2^{-m}$ and, after
normalization, identifies exactly one pair of otherwise independent
clique labels. It therefore has the moment in (b). The union measure is
dominated by the sum of these six positive measures. This gives a valid
positive-spectator price without claiming a unit moment for $E_p/\eta$.
For the density comparison before the Gram step, its actual variation
is $\eta$ and its one-slot marginals are uniform. These are separate uses of the deleted measure $E_p$.

For a uniform alphabet of size $Q$, ordinary H\"older and finite-space
norm comparison give, for every positive uniform-marginal law on $72$ slots,
\[
 |\Lambda((f_v)_v)|\le Q^{\theta-1}
                    \prod_v\|f_v\|_{72/\theta}.
\]
Apply this only to the terms of the selection expansion (\cref{def:price}) supported at the unchanged primes $p<2^{16}$, where $Q=p^{2m}$. An unselected prefix stays in the RP outside
reference and pays no such cost. Complete selection of a clique prime
has the existing conductor price $p^{2ma}$, not $p^{4ma}$.
Bounding the number of primes by the number of integers gives the common
upper bound, for each of $S,A,T$,
\begin{align*}
 &12\cdot2^{16+m(-1+32(a+\theta-1))}
   +12\cdot2^{512+m(-1+1024a)}\\
 &\qquad\le12\left(2^{-2694}+2^{-46}\right)
       <2\cdot10^{-13}<10^{-12}.
\end{align*}
The factor $12$ is $6$ for the six equality events times $2$ for bounding the symmetric and skew parts separately. The first term is also less than $2^{-2600}$.
The new large-prime unit moment leaves only the usual conductor price
in the second term. This proves (d).
\end{proof}

\begin{remark}
The prefix replacement changes no coordinate group, conductor cutoff,
external index variable or square-affine map. The fixed-prefix conductor
argument and the actual-square kernel therefore retain their stated
scope. It changes the local law at the large class-one prefix primes,
and the exceptional deletion can now have both symmetric and skew parts.
Its moment and bracket proofs above account for that change explicitly.
The conclusion is the list of moment and bracket properties of the selected prefix laws assumed in \cref{sec:density-upgrade},
not a stronger moment assertion for arbitrary admissible measures.
\end{remark}
